\documentclass[11pt]{article}
\usepackage[utf8]{inputenc}
\usepackage[french]{babel}
\usepackage{amsmath}
\usepackage{amssymb}
\usepackage{geometry}
\geometry{a4paper, margin=2.5cm}

\title{Analyse des attaques identifiées par Tamarin\\sur le protocole ASKO-OM8464A2}
\author{Garance Frolla, Ely Marthouret, Ewan Decima\\Team: ASKO OM8464A2}
\date{Novembre 2025}

\begin{document}

\maketitle

\section{Introduction}

L'analyse formelle du protocole ASKO-OM8464A2 avec l'outil Tamarin a révélé que deux lemmes de 
sécurité ne peuvent pas être prouvés automatiquement. Ces lemmes concernent les propriétés 
d'authentification et d'accord injectif. Nous présentons dans ce document les traces d'attaque 
générées par Tamarin et expliquons pourquoi ces attaques ne sont pas considérées comme des 
vulnérabilités réelles du protocole dans le contexte de nos spécifications.

\section{Attaque sur le lemme \texttt{aliveness\_B\_to\_A}}

\subsection{Description du lemme}

Le lemme \texttt{aliveness\_B\_to\_A} cherche à vérifier que si l'agent $B$ termine le protocole 
avec l'agent $A$, alors $A$ a également terminé le protocole. Formellement, le lemme s'exprime 
comme suit :    
\[
\forall A, B, k_{AB}, \#j : \text{Finish\_B}(A, B, k_{AB}) @ \#j \implies 
\exists \#i : \text{Finish\_A}(A, B, k_{AB}) @ \#i
\]
où l'on néglige les cas de compromission des clés.

\subsection{Trace d'attaque identifiée par Tamarin}

Tamarin génère une trace d'attaque où l'adversaire exploite le fait que le dernier message du 
protocole peut être intercepté ou perdu. La trace se déroule comme suit. L'agent $A$ initie le 
protocole en générant le nonce $N_A$ et la clé de session $k_{AB}$, puis envoie le premier 
message $\{|\langle A, N_A \rangle|\}_{k_{AB}}$ vers $B$. Ensuite, $A$ contacte le serveur $S$ 
en envoyant $\{|\langle B, \tau, \lambda, k_{AB} \rangle|\}_{k_{AS}}$. Le serveur $S$ transmet 
alors à $B$ le message $\{|\langle A, \tau, \lambda, k_{AB} \rangle|\}_{k_{BS}}$. L'agent $B$ 
reçoit ce message, déchiffre le premier message de $A$, et envoie la confirmation 
$\{|N_A + 1|\}_{k_{AB}}$.

À ce stade, l'adversaire intercepte le quatrième message avant qu'il n'atteigne $A$. Le résultat 
est que $B$ a terminé le protocole avec succès car il a envoyé sa confirmation, mais $A$ n'a 
jamais reçu cette confirmation et reste donc dans un état d'attente. Tamarin considère que $B$ 
a exécuté l'action $\text{Finish\_B}(A, B, k_{AB})$ sans qu'il existe d'action correspondante 
$\text{Finish\_A}(A, B, k_{AB})$.

\subsection{Justification de la validité du protocole}

Cette attaque n'est pas considérée comme une vulnérabilité réelle du protocole ASKO-OM8464A2. 
En effet, nos spécifications précisent explicitement que la validation de l'échange repose sur 
l'envoi du dernier message par $B$. Lorsque $B$ envoie le message $\{|N_A + 1|\}_{k_{AB}}$, il 
démontre qu'il possède la clé de session $k_{AB}$ et qu'il a bien reçu le nonce $N_A$ de $A$. 
Du point de vue de la sémantique du protocole, l'échange est considéré comme valide dès que $B$ 
termine son exécution, indépendamment de la réception effective du message par $A$.

Cette conception reflète un choix architectural où la responsabilité de la validation finale 
incombe au répondant $B$. Dans notre modèle de menace, il est acceptable que $A$ ne reçoive pas 
toujours la confirmation finale, car le protocole garantit que si $B$ termine, alors $A$ a 
nécessairement initié la communication et que la clé $k_{AB}$ a été correctement établie. 
L'important est que $B$ puisse authentifier $A$ via le serveur $S$, ce qui est effectivement 
garanti par le message 3 du protocole.

\section{Attaque sur le lemme \texttt{injective\_agreement\_B\_A}}

\subsection{Description du lemme}

Le lemme \texttt{injective\_agreement\_B\_A} vise à garantir une propriété d'accord injectif 
entre $B$ et $A$. Cette propriété stipule que chaque acceptation de la clé $k_{AB}$ par $B$ 
correspond à une unique initiation par $A$. Formellement :
\[
\forall A, B, t, \#i : \text{Commit\_B}(A, B, \langle \text{`init'}, t \rangle) @ \#i \implies 
\]
\[
\exists \#j : \text{Running\_A}(A, B, \langle \text{`init'}, t \rangle) @ \#j \land \#j < \#i \land
\]
\[
\neg \exists A_2, B_2, \#i_2 : \text{Commit\_B}(A_2, B_2, \langle \text{`init'}, t \rangle) @ \#i_2 \land \#i_2 \neq \#i
\]

\subsection{Trace d'attaque identifiée par Tamarin}

Tamarin génère une trace d'attaque exploitant la possibilité de rejouer les messages du protocole.
Dans un premier temps, une session légitime se déroule entre $A$ et $B$. L'agent $A$ génère la 
clé $k_{AB}$ et le nonce $N_A$, puis les messages du protocole sont échangés normalement jusqu'à
ce que $B$ termine avec succès. L'adversaire enregistre tous les messages échangés durant cette 
session, en particulier les messages $\{|\langle A, N_A \rangle|\}_{k_{AB}}$ et $\{|\langle A, \tau, \lambda, k_{AB} \rangle|\}_{k_{BS}}$.

Par la suite, l'adversaire rejoue ces deux messages vers $B$. Puisque les messages sont toujours 
valides selon la vérification temporelle $t < \tau + \lambda$, l'agent $B$ accepte une deuxième 
fois la clé $k_{AB}$. Du point de vue de Tamarin, il existe maintenant deux exécutions distinctes 
de $\text{Commit\_B}$ avec la même clé $k_{AB}$, alors qu'il n'existe qu'une seule exécution de 
$\text{Running\_A}$. Cela viole la propriété d'injectivité qui exige une correspondance bijective 
entre les acceptations de $B$ et les initiations de $A$.

Cette attaque est rendue possible par l'absence d'un mécanisme liant de manière unique chaque 
réponse de $B$ à une requête spécifique. Le timestamp $\tau$ et la durée de validité $\lambda$ 
ne suffisent pas à empêcher le rejeu durant la fenêtre temporelle $[\tau, \tau + \lambda]$, et 
$B$ ne conserve pas de trace des clés de session déjà utilisées.

\subsection{Justification de la validité du protocole}

Cette attaque par rejeu n'est pas considérée comme une violation de nos spécifications de 
sécurité. En effet, le document de spécifications stipule explicitement que B doit vérifier
que le nonce n'a pas déjà été envoyé. Donc si il reçoit 2 fois le même message il pourra avorter
le protocole.

La propriété d'injectivité stricte n'est donc pas nécessaire dans notre contexte d'utilisation.
Ce qui importe est que chaque session soit correctement authentifiée et que la clé reste secrète 
pendant sa durée d'utilisation effective, ce qui est garanti par les propriétés de fraîcheur et 
de secret que Tamarin a pu vérifier avec succès.

\section{Conclusion}

L'analyse formelle avec Tamarin a permis d'identifier deux traces d'attaque théoriques sur le 
protocole ASKO-OM8464A2. La première concerne l'absence de garantie que l'initiateur $A$ termine 
le protocole lorsque le répondant $B$ le termine, et la seconde porte sur la possibilité de 
rejouer les messages pour violer la propriété d'accord injectif. Toutefois, ces attaques ne 
constituent pas des vulnérabilités dans le cadre de nos spécifications. Le premier cas est 
conforme à notre conception où la validation repose sur l'envoi du dernier message par $B$, 
et le second cas est exclu par notre règle de vérification des nonces.

\end{document}